\section{Introduction}
Lexical URL classifiers need only a URL string, avoiding page fetches, DNS lookups and third-party reputation services during feature extraction. Our two primary papers motivate complementary aspects of this application. S\'anchez-Paniagua~\textit{et al.}~\cite{sanchez2022} compare URL-based detection models, including random forest (RF) and LightGBM, and show why the representation of legitimate login URLs matters. Rashid~\textit{et al.}~\cite{rashid2024} investigate cross-source generalisation and address distribution differences through unsupervised domain adaptation (UDA). Together these studies motivate testing established classifiers beyond their source collection rather than interpreting within-source performance as evidence of general robustness.

\textit{Evaluation gap.} UDA legitimately assumes access to unlabelled target data. We examine a different constraint: no target observations for fitting, feature selection or threshold calibration. Under that constraint, can stability of feature importance across source-domain folds improve transfer over mean importance at the \emph{same} feature budget, and is the reduction in detection cost worth the loss relative to all features? This is a question about the joint evaluation of ranking, budget and transfer, rather than a claim that cross-dataset detection or feature stability has not been studied. URL-level splits can share registered domains; we instead require disjoint domains and evaluate both transfer directions.

\textit{Scope of contribution.} This is an application and empirical benchmark project, not a reproduction of the two primary papers or a new tree-learning algorithm. RF and LightGBM are established classifiers evaluated in~\cite{sanchez2022}; that paper's strongest overall pipeline uses TF-IDF/character n-grams with logistic regression, while~\cite{rashid2024} studies UDA. We do not claim to reproduce or outperform either system. We make three concrete contributions:
\begin{enumerate}
\item A curated, \emph{domain-disjoint} benchmark built from two public sources, EBBU-2017 (also examined in~\cite{sanchez2022}) and HISPAR--Phishstats (HP)~\cite{rashid2024}: registered domains shared by the two sources are removed symmetrically, and every split is group-disjoint by registered domain. The original repositories and commits are retained as dataset provenance.
\item A transparent comparison of two source-only selectors built on the same permutation-importance (PI) evidence (ranking by mean PI, MEAN-PI, or by cross-fold stability, STABLE-PI) at matched feature budgets, for RF and LightGBM, in both transfer directions, with nested selection, thresholds locked on source validation data and domain-cluster bootstrap intervals.
\item An empirical assessment of transfer--cost trade-offs, including negative comparisons against ALL features, with parsing+extraction, prediction-only and end-to-end latency measured separately. SHAP is a secondary diagnostic with the validation limitations reported below.
\end{enumerate}
The claimed distinction is this controlled evaluation and its bounded empirical insight; global methodological novelty is not established. Related stability-based work already exists~\cite{kanon2026,misiek2026}. Target identities are used to curate the domain-disjoint cohort, and VAL URLs from both sources form the timing pool; thus \emph{source-only} describes learning and calibration, not target-blind benchmark acquisition. The experimental protocol was frozen before evaluation; this editorial positioning was clarified afterwards without changing the planned contrasts or results.
